\subsection{Kinematic distributions after the 2021 preselection}
\label{sec:preselection-kinematics}
The five groups in Figures~\ref{fig:presel_vertex}--\ref{fig:scatter_chi2}
show the reconstructed quantities used to inspect the 2021 loose preselection.
Figure~\ref{fig:presel_vertex} displays the mass spectrum and vertex-fit quality;
Figures~\ref{fig:presel_ele_kin} and~\ref{fig:presel_pos_kin} show the electron
and positron momenta and impact parameters. The two-dimensional displays in
Figures~\ref{fig:scatter_mom} and~\ref{fig:scatter_chi2} make their mass dependence
visible and distinguish broad kinematic structure from the narrow simulated
mass hypotheses. The signal samples are scaled for display, so these figures
compare reconstruction shapes rather than absolute yields or efficiencies.
The vertex-fit $\chi^2$ remains a diagnostic: no additional vertex-fit $\chi^2$
requirement is applied to the regenerated 2021 GPR input. These archived plots
document the preselection context and do not constitute a new event selection
or a validation of the final likelihood.

\begin{figure}[htbp]
  \centering
  \begin{subfigure}[t]{0.48\textwidth}
    \includegraphics[width=\textwidth]{preselection_plots/signal_overlay_loose_loose_vtx_mass.pdf}
    \caption{Vertex invariant mass.}
  \end{subfigure}\hfill
  \begin{subfigure}[t]{0.48\textwidth}
    \includegraphics[width=\textwidth]{preselection_plots/signal_overlay_loose_loose_vtx_chi2.pdf}
    \caption{Vertex-fit $\chi^2$.}
  \end{subfigure}
  \caption{Reconstructed-vertex distributions after the loose preselection,
    comparing data (points) with the signal MC (filled histogram). The vertex-fit
    $\chi^2$ is shown as a diagnostic; no vertex-fit $\chi^2$ requirement is applied to
    the regenerated 2021 GPR input.}
  \label{fig:presel_vertex}
\end{figure}

% ------------------------------------------------------------
%  1. Electron track kinematics
% ------------------------------------------------------------
% The following figures compare the loose-preselected data (points) with the
% signal Monte Carlo (histogram), normalised to unit area, for the kinematic
% variables entering the preselection. Figure~\ref{fig:presel_ele_kin} shows the
% electron-track quantities: the total and longitudinal momentum together with
% the transverse and longitudinal impact parameters. Good agreement between the
% shapes of the reconstructed distributions is expected in the bulk of phase
% space, while the signal is characterised by its harder momentum spectrum and
% displaced-vertex tails.
\begin{figure}[htbp]
  \centering
  \begin{subfigure}[t]{0.48\textwidth}
    \includegraphics[width=\textwidth]{preselection_plots/signal_overlay_loose_loose_ele_p.pdf}
    \caption{Electron momentum $p$.}
  \end{subfigure}\hfill
  \begin{subfigure}[t]{0.48\textwidth}
    \includegraphics[width=\textwidth]{preselection_plots/signal_overlay_loose_loose_ele_pz.pdf}
    \caption{Electron longitudinal momentum $p_z$.}
  \end{subfigure}

  \vspace{0.5em}
  \begin{subfigure}[t]{0.48\textwidth}
    \includegraphics[width=\textwidth]{preselection_plots/signal_overlay_loose_loose_ele_d0.pdf}
    \caption{Electron transverse impact parameter $d_0$.}
  \end{subfigure}\hfill
  \begin{subfigure}[t]{0.48\textwidth}
    \includegraphics[width=\textwidth]{preselection_plots/signal_overlay_loose_loose_ele_z0.pdf}
    \caption{Electron longitudinal impact parameter $z_0$.}
  \end{subfigure}
  \caption{Electron-track kinematic distributions after the loose preselection,
    comparing data (points) with the signal MC (filled histogram).}
  \label{fig:presel_ele_kin}
\end{figure}

% ------------------------------------------------------------
%  2. Positron track kinematics
% ------------------------------------------------------------
% Figure~\ref{fig:presel_pos_kin} shows the corresponding positron-track
% quantities. As for the electron, the momentum and impact-parameter
% distributions are used to validate the modelling of the reconstructed final
% state prior to any signal-region selection.
\begin{figure}[htbp]
  \centering
  \begin{subfigure}[t]{0.48\textwidth}
    \includegraphics[width=\textwidth]{preselection_plots/signal_overlay_loose_loose_pos_p.pdf}
    \caption{Positron momentum $p$.}
  \end{subfigure}\hfill
  \begin{subfigure}[t]{0.48\textwidth}
    \includegraphics[width=\textwidth]{preselection_plots/signal_overlay_loose_loose_pos_pz.pdf}
    \caption{Positron longitudinal momentum $p_z$.}
  \end{subfigure}

  \vspace{0.5em}
  \begin{subfigure}[t]{0.48\textwidth}
    \includegraphics[width=\textwidth]{preselection_plots/signal_overlay_loose_loose_pos_d0.pdf}
    \caption{Positron transverse impact parameter $d_0$.}
  \end{subfigure}\hfill
  \begin{subfigure}[t]{0.48\textwidth}
    \includegraphics[width=\textwidth]{preselection_plots/signal_overlay_loose_loose_pos_z0.pdf}
    \caption{Positron longitudinal impact parameter $z_0$.}
  \end{subfigure}
  \caption{Positron-track kinematic distributions after the loose preselection,
    comparing data (points) with the signal MC (filled histogram).}
  \label{fig:presel_pos_kin}
\end{figure}



% ------------------------------------------------------------
%  5. Two-dimensional mass correlations
% ------------------------------------------------------------
% The two-dimensional distributions of the reconstructed invariant mass versus
% the preselection variables are shown for data (left column) and signal MC
% (right column) in the following figures. These correlations are used to check
% that the mass reconstruction is unbiased across the relevant phase space and
% that no artificial structure is introduced by the selection.
% Figure~\ref{fig:scatter_mom} shows the mass-versus-momentum correlations.
\begin{figure}[htbp]
  \centering
  \begin{subfigure}[t]{0.48\textwidth}
    \includegraphics[width=\textwidth]{preselection_plots/xy_scatter_xy_mass_vs_ele_p_data_loose.pdf}
    \caption{Mass vs.\ electron $p$ (data).}
  \end{subfigure}\hfill
  \begin{subfigure}[t]{0.48\textwidth}
    \includegraphics[width=\textwidth]{preselection_plots/xy_scatter_xy_mass_vs_ele_p_signal_loose.pdf}
    \caption{Mass vs.\ electron $p$ (signal).}
  \end{subfigure}

  \vspace{0.5em}
  \begin{subfigure}[t]{0.48\textwidth}
    \includegraphics[width=\textwidth]{preselection_plots/xy_scatter_xy_mass_vs_pos_p_data_loose.pdf}
    \caption{Mass vs.\ positron $p$ (data).}
  \end{subfigure}\hfill
  \begin{subfigure}[t]{0.48\textwidth}
    \includegraphics[width=\textwidth]{preselection_plots/xy_scatter_xy_mass_vs_pos_p_signal_loose.pdf}
    \caption{Mass vs.\ positron $p$ (signal).}
  \end{subfigure}

  \caption{Reconstructed invariant mass versus momentum observables, for data
    (left) and signal MC (right), after the loose preselection.}
  \label{fig:scatter_mom}
\end{figure}

% Figure~\ref{fig:scatter_chi2} shows the mass-versus-track/vertex-quality
% correlations.
\begin{figure}[htbp]
  \centering
  \begin{subfigure}[t]{0.48\textwidth}
    \includegraphics[width=\textwidth]{preselection_plots/xy_scatter_xy_mass_vs_ele_chi2_data_loose.pdf}
    \caption{Mass vs.\ electron $\chi^2$ (data).}
  \end{subfigure}\hfill
  \begin{subfigure}[t]{0.48\textwidth}
    \includegraphics[width=\textwidth]{preselection_plots/xy_scatter_xy_mass_vs_ele_chi2_signal_loose.pdf}
    \caption{Mass vs.\ electron $\chi^2$ (signal).}
  \end{subfigure}

  \vspace{0.5em}
  \begin{subfigure}[t]{0.48\textwidth}
    \includegraphics[width=\textwidth]{preselection_plots/xy_scatter_xy_mass_vs_pos_chi2_data_loose.pdf}
    \caption{Mass vs.\ positron $\chi^2$ (data).}
  \end{subfigure}\hfill
  \begin{subfigure}[t]{0.48\textwidth}
    \includegraphics[width=\textwidth]{preselection_plots/xy_scatter_xy_mass_vs_pos_chi2_signal_loose.pdf}
    \caption{Mass vs.\ positron $\chi^2$ (signal).}
  \end{subfigure}

  \vspace{0.5em}
  \begin{subfigure}[t]{0.48\textwidth}
    \includegraphics[width=\textwidth]{preselection_plots/xy_scatter_xy_mass_vs_vtx_chi2_data_loose.pdf}
    \caption{Mass vs.\ vertex $\chi^2$ (data).}
  \end{subfigure}\hfill
  \begin{subfigure}[t]{0.48\textwidth}
    \includegraphics[width=\textwidth]{preselection_plots/xy_scatter_xy_mass_vs_vtx_chi2_signal_loose.pdf}
    \caption{Mass vs.\ vertex $\chi^2$ (signal).}
  \end{subfigure}
  \caption{Reconstructed invariant mass versus track- and vertex-fit quality,
    for data (left) and signal MC (right), after the loose preselection.}
  \label{fig:scatter_chi2}
\end{figure}


\FloatBarrier
